% Purpose: what the results mean physically (why structure suffices, local vacancy chemistry), how they relate to prior work, and the scope (home of the statement that no experimental formation energies were used). One "We propose" per physical reading, with the measurement that motivates it and the one that would test it. Register: docs/prose_style.md. Every number is a macro of paper/generated/numbers.tex.
\section{Discussion}
\label{sec:discussion}

\subsection{Why Structure Suffices}
\label{sec:disc_principal}

From \numNetNStar{} training oxides on, the electronic structure of the pristine host lowers the site error by at most \numNetHeadlineAUpper{}~eV, inside the tolerance, while with \numCfgBudgetsHostsItemFive{} its advantage is not yet bounded below $\delta$ (Section~\ref{sec:learning_curves}). The question is why explicit electronic information stops paying for itself at a few hundred oxides.

Part of the answer is a fact about the descriptors. The band gap, the oxide formation energy, the dielectric response, the orbital character of the band edges, the Bader and Born effective charges and the centre of the oxygen $p$ band are all outputs of one DFT calculation on the pristine cell (Section~\ref{sec:data}), so each is fixed by the composition and atomic positions that the graph also receives. A map from the pristine structure to $E_\mathrm{f}$ can therefore absorb them, and the learning curve measures the sample complexity of doing so: how many oxides an encoder of the coordination environment needs before it reproduces the part of this information that matters for the vacancy. The coordination of the oxygen, the strength of its M--O bonds and the stability of the lattice set how much the broken bonds cost and where the two electrons can go, and all three are written in the structure.

We propose that what the descriptors still contribute is mainly host thermodynamics, which a structure-only network learns readily. Three measurements point there. The host carries most of the variance of $E_\mathrm{f}$ ($\mathrm{ICC}_\mathrm{host} = \numVarianceDecompositionHostIcc{}$), and B0, linear in the oxide formation energy and the band gap, already reaches \numNetSixHundredFiftyFourMaeBZero{}~eV. The descriptor predictor recovers the oxide formation energy best of all descriptors, with a coefficient of determination of \numStagingVTwoPOneDescriptorsSixHundredFiftyFourPerDescriptorDescFormationEnergyRTwoMean{}, and the host-electronic class better than the site-electronic class (Section~\ref{sec:staging}). And the forest on structural and compositional descriptors alone, at \numBaselinesSummaryOuterFullBudgetRfStructuralMaeMean{}~eV, is more accurate than the forest on electronic descriptors alone, at \numBaselinesSummaryOuterFullBudgetRfElectronicMaeMean{}~eV, so what places a host can be read without an electronic-structure calculation; an untrained encoder of the same coordination environments already exposes much of the descriptor variance (Appendix~\ref{app:probe}). The hypothesis makes two predictions. Given to D separately, the host-electronic class should carry nearly all of $A$ at every training size. For charged vacancies and shallow donors, whose energies depend on the position of the band edges directly rather than through the local bonding, explicit electronic information should keep its advantage to far larger training sets.

Staging costs accuracy for a related reason: the descriptors it passes on carry their largest prediction errors in the site-level charges, which describe the local ionic environment of the removed oxygen and which the pristine structure fixes in principle but a network trained on a few hundred oxides does not yet reproduce (the Bader charge reaches a coefficient of determination of \numStagingVTwoPOneDescriptorsSixHundredFiftyFourPerDescriptorDescBaderChargeRTwoMean{}). Feeding D the true descriptors with noise added at the per-descriptor error of P1, class by class, would locate that cost.

\subsection{Local Vacancy Chemistry}
\label{sec:disc_site}

Within one oxide the vacancy energy changes from site to site far less than between oxides, and neither network captures most of that variation: with \numCfgBudgetsMax{} training hosts the within-host residual MAE of S, \numCOneCTwoVTwoMaeSixHundredFiftyFourSWithinHostMaeMean{}~eV, removes about a fifth of the zero-skill value of \numHostMeanOracleAggregateWithinHostMaeMean{}~eV (Section~\ref{sec:within_host}). We propose that within a host the vacancy energy is ordered by the local bonding and by the relaxation of the neighbouring cations rather than by the charge of the oxygen site. That reading explains why the Bader charge ranks the sites of a host with a correlation of only \numWithinHostBreakdownSpearmanReferenceBaderChargeMean{} and why pristine-site descriptors give D no within-host advantage; a breakdown of the within-host residuals by coordination number and M--O bond length would test it.

The cation $d$ shell sets the size of the spread: the zero-skill value is \numWithinHostBreakdownGroupsDClassDZeroVTwoSZeroSkillMaePoint{}~eV in \dzero{} hosts and \numWithinHostBreakdownGroupsDClassDTenVTwoSZeroSkillMaePoint{}~eV in hosts with a filled $d$ shell. In closed-shell \dzero{} hosts the two electrons stay at the vacancy and sites differ mostly in the number of cation--oxygen bonds broken, which the neighbour list of the pristine structure encodes directly. With filled-$d$-shell cations such as Zn, Ga, Cd and Sn the M--O bonds are more covalent, and the cost depends on which soft, polarisable cations are bonded to the site and on how the vacancy electrons hybridise with their $s$ states. The screened pairs follow this picture. In Ga$_2$O$_3$ the four-coordinate site, with the most bonds to break, is the highest in DFT and in the prediction, whereas in ZrO$_2$ the four-coordinate site, with the longer mean Zr--O distance, lies below a three-coordinate site in DFT, and the network inverts the two (Section~\ref{sec:screening}): there longer, weaker bonds outweigh their number, a distinction the network has not learned from \numNetScreenBudget{} oxides.

\subsection{Relation to Prior Work}
\label{sec:disc_prior}

Kumagai et al.\ built the database and showed that random forests on descriptors of the pristine host and site predict the neutral vacancy energy, with a learning curve over the number of training oxides \citep{Kumagai2021insights}; our RF-Kumagai reproduces their error to within \numBaselinesSummaryRfKumagaiVsPublishedDifferenceEV{}~eV. Kiyohara et al.\ showed that a graph network on the crystal structure alone reaches \numNetPublishedMae{}~eV for the neutral vacancy \citep{Kiyohara2025machinelearning}; the structure-only accuracy is theirs, and our network matches it on formula-grouped splits. Witman et al.\ predicted site-resolved vacancy formation enthalpies of oxides from the ideal crystal structure with defect graph networks \citep{Witman2023defect}, and Deml et al.\ identified the oxide formation enthalpy and the band gap as the controls of the vacancy energy \citep{Deml2014oxide,Deml2015intrinsic}, which B0 encodes. Linton, Singh and Aidhy predicted Bader charges and then vacancy energies in face-centred-cubic high-entropy alloys \citep{Linton2026framework}; our staged model applies that serial design to oxides with every electronic descriptor of the database.

What is new is the training size at which explicit pristine-host electronic information stops lowering the error of a structure-only network by more than a chemically chosen tolerance, measured for both inputs on the same held-out oxides with redrawn training sets and equal tuning, together with the site-level, staged and relaxed-geometry tests of the structure-only route.

\subsection{Scope}
\label{sec:disc_scope}

The results describe deep, localised neutral vacancies in ordered non-magnetic oxides computed with one functional and one workflow (Section~\ref{sec:data}). The charged vacancies and the shallow-donor entries of the same database are the next test. The formation energy of $V_\mathrm{O}^{\bullet\bullet}$ depends on the electron chemical potential and so on the band-edge positions \citep{Freysoldt2014firstprinciples}, and the energy of a shallow donor is set by the conduction band into which its electrons delocalise \citep{Lany2005anion}; Section~\ref{sec:disc_principal} predicts that the descriptors keep their advantage there.

The domain leaves out many of the oxides in which vacancies are put to use. Among the screened candidates, In$_2$O$_3$, CeO$_2$ and WO$_3$ have no host in the release and SrTiO$_3$ and BaTiO$_3$ leave it as shallow donors (Section~\ref{sec:screening}), although they are among the most studied conducting, reducible and electrochromic oxides \citep{Ellmer2012past,Esch2005electron,Granqvist2000electrochromic}. Within the domain, the two energy groups of the screened sites follow the host band gap, with a rank correlation of \numScreeningDemoDescriptiveSpearmanDftVsBandGap{} across the \numNetScreenSites{} sites: in wide-gap, strongly bound oxides the broken bonds are costly and the electrons must occupy a level deep in the gap, whereas the cation-$s$ conduction bands of ZnO, SnO$_2$ and Ga$_2$O$_3$ take them up at lower cost. For cation families absent from training the advantage is not bounded below $\delta$ (Section~\ref{sec:learning_curves}).

No experimental formation energies were used to validate the predictions; every comparison is with the DFT energies of the release.
